\section{Industrial Model-Based Tools and Ecosystems}
\label{sec:sota:tools}

The automotive industry has long relied on a specific set of model-based tools. This section
evaluates these dominant and emerging tools against the challenges of the Software-Defined Vehicle
(\sdv), demonstrating a clear gap in the current state of the art that \GRust aims to fill.

\subsection{The Legacy: Pull-Based, C-Centric Toolchains}
\label{sec:sota:tools:legacy}

The established toolchains for automotive design, while proven and effective, are built on a
paradigm that creates significant architectural friction with the emerging needs of the \sdv.

\paragraph*{Simulink}
As the de facto industry standard, \Simulink~\cite{simulink} offers a powerful graphical
environment for modeling and simulation. However, its execution model, based on discrete-time
solvers and periodic sample times, is the direct source of the inefficiencies identified in
\Cref{sec:intro:gap}. This pull-based, periodic approach is ill-suited for efficiently modeling the
event-driven reality of an \sdv. Furthermore, its generation of standard \Ci
code~\cite{simulink_codegen,simulink_misra} forgoes the memory and concurrency safety offered by
modern languages. While highly optimized for legacy architectures, it is less adapted to the dynamic
software environment of the \sdv.

\paragraph*{Scade Suite}
For safety-critical systems, \Scade~\cite{DBLP:conf/fdl/ColacoPPP18} represents the pinnacle of the
synchronous, model-based approach, offering formal guarantees and certifiable code generation based
on \Lustre~\cite{DBLP:conf/popl/CaspiPHP87}. However, its design philosophy is deeply rooted in the
single-threaded, synchronous world and lacks native constructs for asynchronous dataflow. This poses
significant challenges for a seamless and natural integration with event-driven runtimes like
\texttt{tokio}. While formally sound, \Scade presents an architectural mismatch with the emerging,
asynchronous automotive software ecosystem.

\subsection{The Emerging Alternative: Rust-Centric Stacks}
\label{sec:sota:tools:rust_stacks}

In response to the limitations of legacy tools, new \Rust-centric stacks have emerged that embrace
an asynchronous, event-driven model. However, they introduce a new abstraction gap.

\paragraph*{Veecle}
This project provides a modular runtime and high-level APIs for building asynchronous automotive
applications in \Rust. It successfully abstracts away middleware and hardware details, but it is
primarily a runtime infrastructure, not a model-based design tool. It provides the framework of
execution but lacks the high-level, domain-specific modeling abstractions familiar to control
engineers, leaving them to implement complex logic using low-level programming primitives.

\paragraph*{Pictorus}
This tool addresses the modeling gap by offering a visual, \Simulink-like environment that generates
\Rust code. It successfully bridges the language gap but fails to bridge the architectural one. Its
underlying execution model remains largely sequential, lacking the native, first-class support for
event-driven dataflow and asynchronous interactions that are essential for building efficient \sdv
software.

\subsection{Synthesis: The Unfilled Gap}
\label{sec:sota:tools:synthesis}

The current landscape presents a clear dichotomy. On one side, legacy tools like \Simulink and
\Scade offer mature, high-level modeling abstractions but are tied to an inefficient, pull-based
execution model. On the other side, emerging \Rust-based stacks embrace the correct event-driven
architecture but fail to provide the high-level, formally-grounded modeling abstractions needed for
complex control systems.

Neither of these worlds successfully marries the rigor of synchronous, model-based design with the
safety and asynchronous efficiency of the modern \Rust-based \sdv platform. This unfilled gap is
precisely where \GRust is positioned. It is designed to provide the formal, declarative modeling of
a synchronous language while directly generating safe asynchronous \Rust code, bridging the gap that
current industrial and emerging tools leave open.

\section{Synchronous Languages and Academic Toolsets}
\label{sec:sota:synchronous}

The synchronous paradigm, which models computation as occurring in a sequence of discrete,
instantaneous reactions, has given rise to a rich family of academic languages. These languages
provide a formal basis for designing critical systems with strong guarantees of correctness and
predictability. This section surveys influential synchronous languages, analyzing their distinct
approaches to managing time, concurrency, and program structure.

\subsection{Lustre: A Declarative Synchronous Language}
\label{sec:sota:synchronous:lustre}

As introduced in \Cref{sec:intro:dsl:dataflow}, a dataflow model describes a system as a network of
components that process data flows. \Lustre is a \emph{declarative} dataflow language where programs
are defined as pure functions over these flows, which are called \emph{streams}~\cite{DBLP:conf/popl/CaspiPHP87}.

In \Lustre, a stream is an infinite sequence of values aligned on a discrete \emph{logical clock}.
All streams are \emph{synchronous}, meaning their values are conceptually produced at the same
logical instants. A program, called a \emph{node}, is defined by a set of equations that specify how
output streams are computed from input streams. The \texttt{pre} operator provides access to the
value of a stream at the previous instant, enabling stateful computations.

The central concept in \Lustre is the \emph{clock} of a stream, which defines the sequence of
instants at which it holds a \emph{concrete value}. At other instants, the stream is considered
absent and carries a special \emph{bottom value}, denoted by $\bot$. The base clock represents the
fastest rate of reaction. Slower streams can be derived using the \texttt{when} operator, which
samples a stream based on a boolean condition. The compiler uses a \emph{clock calculus} to
statically verify that operations only combine streams that are present at the same instants. This
formal verification guarantees bounded memory usage, determinism, and a predictable reaction time
(\wcet), making \Lustre a foundational language for safety-critical systems, as demonstrated by its
industrial successor, \Scade~\cite{DBLP:conf/fdl/ColacoPPP18}.

\begin{multicols}{2}[The \texttt{EXAMPLE} program below illustrates these features.]
  \begin{lstlisting}[basicstyle=\ttfamily\footnotesize\color{BlackText}]
  node EXAMPLE(I: int) returns
    (CLK: bool; O: int when CLK)
  var V: int;
  let
    CLK = true -> not pre(CLK);
    V = 0 -> pre(I);
    O = (V * 2) when CLK;
  tel
\end{lstlisting}

  \newcolumn

  \hfill

  \begin{table}[H]
    \centering
    \begin{tabular}{l|cccccc}
      \texttt{I}   & 1             & 2              & 3             & 4              & \dots \\
      \texttt{CLK} & \texttt{true} & \texttt{false} & \texttt{true} & \texttt{false} & \dots \\
      \texttt{V}   & 0             & 1              & 2             & 3              & \dots \\
      \texttt{O}   & 0             & $\bot$         & 4             & $\bot$         & \dots \\
    \end{tabular}
    \caption{Execution of the \texttt{EXAMPLE} program.}
    \label{tab:sota:synchronous:lustre:exec}
  \end{table}
\end{multicols}

\Cref{tab:sota:synchronous:lustre:exec} shows an execution trace of the program. The node computes
two outputs, \texttt{CLK} and \texttt{O}, from the input \texttt{I}. The equation
\texttt{CLK = true -> not pre(CLK)} defines a boolean stream \texttt{CLK} that is \texttt{true} at the
first instant and alternates at subsequent instants. The stream \texttt{V} is initialized to
\texttt{0} and then takes the previous value of \texttt{I}. The output \texttt{O} is defined on the
clock of \texttt{CLK}: it has a concrete value only when \texttt{CLK} is \texttt{true}, at which point
its value is \texttt{V * 2}.

\paragraph*{Comparison with GRust}

\GRust, and specifically its \GRSync component sub-language, is heavily inspired by \Lustre. It mirrors
the declarative, equation-based style and uses a \texttt{last} operator for state management, which
is analogous to the \texttt{pre} operator in \Lustre.

The key difference lies in the handling of clocks and concurrency. \Lustre employs a formal clock
calculus to manage multi-rate behavior within a single, unified synchronous model. \GRSync
simplifies this by omitting user-defined clocks. Instead, multi-rate computations are handled by
built-in constructs, such as \texttt{when} and \texttt{match} equations. While less expressive than
a full clock calculus, this approach is simpler to use. A \GRSync component is thus a single-clocked
synchronous island, which is asynchronously triggered in a push manner by a \GRFRP service,
following the GALS architecture described in the introduction.

This design choice acknowledges the complexity of a full clock calculus. Indeed, modern extensions
to \Lustre have introduced higher-level constructs such as \texttt{automaton} and \texttt{switch}.
While built upon the clock calculus, these features abstract away direct clock manipulation,
offering a more structured and accessible way to model state-based behavior. This evolution in
\Lustre itself suggests a trend toward providing abstractions that are easier to use than the
underlying formal mechanisms.

\subsection{Signal: A Polychronous Relational Language}
\label{sec:sota:synchronous:signal}

Like \Lustre, \Signal is a declarative, dataflow synchronous language, but it follows a
\emph{relational} approach where a program is a system of equations over streams, also called
\emph{signals}~\cite{DBLP:journals/scp/BenvenisteGJ91}.

The key innovation of \Signal is a sophisticated \emph{clock calculus} that treats clocks as
first-class citizens. This allows a programmer to declare clock variables (of type \texttt{event}),
extract the clock from a signal (\eg \texttt{e = when s}), and define explicit equations over these
clocks (\eg \verb|o ^= a ^+ b| to specify that the clock of \texttt{o} is the union of those of
\texttt{a} and \texttt{b}). By making timing relationships a part of the program specification
itself, \Signal provides a powerful formalism for describing complex, multi-rate, and distributed
systems.

This relational model enables \emph{polychrony}~\cite{DBLP:journals/jcsc/GuernicTL03}, meaning a
program can be specified with multiple, independent base clocks. The compiler's task is to prove the
logical and temporal coherence of the system by solving the global set of clock constraints. In
contrast, \Lustre is \emph{endochronous}, assuming a single base clock. While polychrony is highly
expressive, it introduces significant compilation challenges. The underlying constraint-solving
problem is NP-hard, meaning the compiler cannot guarantee a solution for all valid programs and may
reject some correct specifications for the sake of efficiency.

\begin{multicols}{2}[The \texttt{SAMPLER} process below illustrates the relational model.]
  \begin{lstlisting}[basicstyle=\ttfamily\footnotesize\color{BlackText}]
  process SAMPLER =
    (? integer A, B, C;
     ! integer O;)
    (| CLK ^= (when A) ^+ (when B);
     | O   := C when CLK;
     |)
  where
        event CLK
  end
\end{lstlisting}

  \newcolumn

  \hfill

  \begin{table}[H]
    \centering
    \begin{tabular}{l|cccccc}
      \texttt{A}   & 10     & $\bot$ & 12  & $\bot$ & \dots \\
      \texttt{B}   & $\bot$ & 20     & 21  & $\bot$ & \dots \\
      \texttt{C}   & 100    & 101    & 102 & 103    & \dots \\
      \texttt{CLK} & t      & t      & t   & $\bot$ & \dots \\
      \texttt{O}   & 100    & 101    & 102 & $\bot$ & \dots \\
    \end{tabular}
    \caption{Execution of the \texttt{SAMPLER} process.}
    \label{tab:sota:synchronous:signal:exec}
  \end{table}
\end{multicols}

The example defines a process that samples a dense signal \texttt{C} based on an event derived from
two sparse signals, \texttt{A} and \texttt{B}. It explicitly declares a local event, \texttt{CLK},
and defines it with the equation \verb|CLK ^= (when a) ^+ (when b)|. This constrains \texttt{CLK} to
be present at the union of the instants where \texttt{A} or \texttt{B} are present.

The output \texttt{O} is then defined by the equation \texttt{O := c when CLK}, meaning \texttt{O}
takes the value of \texttt{C} only at the instants defined by \texttt{CLK}. As shown in the trace in
\Cref{tab:sota:synchronous:signal:exec}, the compiler synthesizes the clock for \texttt{CLK} from
\texttt{A} and \texttt{B}. The output \texttt{O} is then produced only when \texttt{CLK} is present,
demonstrating how value and timing constraints are solved together.

\paragraph*{Comparison with GRust}

\Signal uses its powerful but complex clock calculus to formally reason about all timing relations
across the entire system. In contrast, \GRust adopts a more constrained GALS model. It isolates
synchronous logic into single-clocked \GRSync components and manages the asynchronous interactions
between them at the \GRFRP service layer. This approach deliberately avoids the need for a global
clock constraint solver, trading the expressive power of polychrony for simpler compilation and a
clear separation of concerns between local computation and global communication.

\subsection{The Esterel Family: Imperative Synchrony}
\label{sec:sota:synchronous:esterel}

In contrast to the declarative style of \Lustre, \Esterel is an imperative synchronous language
designed for programming complex control logic~\cite{DBLP:journals/scp/BerryG92}. It provides
high-level control structures, including sequencing, explicit concurrency (\via the construct
\texttt{||}), and preemption mechanisms. Statements like \texttt{await}, \texttt{emit}, and
\texttt{loop} allow programmers to specify intricate reactive behaviors in a style closer to
traditional imperative programming.

\begin{multicols}{2}[The \texttt{ABRO} module below illustrates the imperative style of \Esterel.]
  \begin{lstlisting}[basicstyle=\ttfamily\footnotesize\color{BlackText}]
  module ABRO:
    input A, B, R;
    output O;
    loop
      [ await A || await B ];
      emit O
    each R
  end module
\end{lstlisting}

  \newcolumn

  \hfill

  \begin{table}[H]
    \centering
    \begin{tabular}{l|cccccc}
      \texttt{A} & \texttt{()} & $\bot$      & $\bot$      & \texttt{()} & $\bot$      & \dots \\
      \texttt{B} & $\bot$      & \texttt{()} & $\bot$      & $\bot$      & \texttt{()} & \dots \\
      \texttt{R} & $\bot$      & $\bot$      & \texttt{()} & $\bot$      & \texttt{()} & \dots \\
      \texttt{O} & $\bot$      & \texttt{()} & $\bot$      & $\bot$      & $\bot$      & \dots \\
    \end{tabular}
    \caption[An execution trace of the \texttt{ABRO} module.]%
    {An execution trace of the \texttt{ABRO} module. \texttt{()} denotes the presence of a signal,
      $\bot$ its absence.}
    \label{tab:sota:synchronous:esterel:exec}
  \end{table}
\end{multicols}

The core statement, \texttt{[ await A || await B ]}, uses the parallel operator \texttt{||} to wait
until both input signals \texttt{A} and \texttt{B} have occurred. It initiates two concurrent
threads of control; one for each \texttt{await}. The statement terminates at the precise instant the
\emph{second} of these signals is received. The sequencing operator \texttt{;} is
\emph{instantaneous}, it ensures that the \texttt{emit O} statement is executed in the very same
instant that the preceding parallel statement terminates. Thus, \texttt{O} is emitted as soon as
both \texttt{A} and \texttt{B} have been seen. The entire behavior is wrapped in a
\texttt{loop...each R} construct. This statement implements a strong preemption: if signal
\texttt{R} is present, the body of the loop is immediately aborted and restarted, discarding its
current state. This gives the reset signal \texttt{R} priority over the internal logic.

The trace in \Cref{tab:sota:synchronous:esterel:exec} reflects this behavior. At the second instant,
\texttt{O} is emitted because \texttt{B} arrives, fulfilling the condition set by the parallel
statement (since \texttt{A} was seen in the first instant). In the third instant, \texttt{R} resets
the module. At he fifth instant, although \texttt{B} occurs (which would have completed the wait
started in the fourth instant), the simultaneous presence of \texttt{R} aborts the computation
\emph{before} \texttt{O} can be emitted. The module is reset, which results in \texttt{O} being
absent in that instant.

The semantics of \Esterel is formally defined by compiling programs into finite state machines or,
more commonly, logic circuits~\cite{DBLP:books/daglib/0023221}. This compilation process resolves
all concurrency and signal dependencies at compile time, producing a deterministic, single-threaded
execution function. \Esterel is then suited for hardware synthesis and time-critical control
applications.

\smallskip

The ideas in \Esterel have been influential, inspiring derivatives that bring its rigorous
concurrency model to other ecosystems.

\paragraph*{ReactiveML}
\textsc{ReactiveML}~\cite{DBLP:conf/ppdp/MandelP05} transposes the synchronous-reactive model of
\Esterel into the functional programming language \Caml. It is implemented as a compiler extension
that translates a dedicated syntax into standard \Caml code. In \textsc{ReactiveML}, data is
manipulated through pure functions, leveraging the type safety and composability inherent to the
functional paradigm. The language retains the core principles of \Esterel, including instantaneous
broadcast of signals, deterministic concurrency, and strong preemption. Programs are structured as a
set of parallel processes that communicate through signals. At each logical instant, processes
execute until they reach a statement that explicitly consumes time, such as \texttt{pause} or
\texttt{await}, ensuring the system evolves deterministically.

\begin{multicols}{2}[The \texttt{EDGE} process below illustrates \textsc{ReactiveML}.]
  \begin{lstlisting}[basicstyle=\ttfamily\footnotesize\color{BlackText}]
  let process EDGE I O =
    loop
      present I
      then pause
      else (await immediate I;
            emit O)
    end
\end{lstlisting}

  \newcolumn

  \hfill

  \begin{table}[H]
    \centering
    \begin{tabular}{l|ccccccc}
      \texttt{I} & $\bot$ & \texttt{()} & \texttt{()} & $\bot$ & \texttt{()} & $\bot$ & \dots \\
      \texttt{O} & $\bot$ & \texttt{()} & $\bot$      & $\bot$ & \texttt{()} & $\bot$ & \dots \\
    \end{tabular}
    \caption[An execution trace of the \texttt{EDGE} process.]%
    {An execution trace of the \texttt{EDGE} process. \texttt{()} denotes the presence of a signal,
      $\bot$ its absence.}
    \label{tab:sota:synchronous:reactiveml:exec}
  \end{table}
\end{multicols}

The process is designed to detect a rising edge: it emits the output signal \texttt{O} only in the
first instant that the input signal \texttt{I} becomes present after an absence. The behavior is
defined by a loop containing a \texttt{present} statement, which checks for the presence of
\texttt{I} in the current instant. If the signal is present, the \texttt{then pause} branch is
taken; the \texttt{pause} statement consumes a single logical instant, waiting for the next tick
without taking further action. This prevents \texttt{O} from being emitted repeatedly while
\texttt{I} remains high. Conversely, if \texttt{I} is absent, the \texttt{else} branch is executed.
Here, the process executes \texttt{await immediate I}, suspending execution until \texttt{I}
appears. The instantaneous sequencing operator (\texttt{;}) then ensures that \texttt{emit O} is
executed in the very same instant that the wait is fulfilled.

The execution trace in \Cref{tab:sota:synchronous:reactiveml:exec} demonstrates this logic. In the
first instant, \texttt{I} is absent, so the process begins to await its arrival. In the second
instant, \texttt{I} appears, the \texttt{await} terminates, and \texttt{O} is emitted. In the third
instant, \texttt{I} is still present, so the \texttt{pause} statement is executed and \texttt{O} is
absent.

\paragraph*{HipHop.js}
\textsc{HipHop.js}~\cite{DBLP:conf/pldi/BerryS20} adapts the synchronous-reactive model to
\JavaScript, aiming to solve the challenges of coordinating complex, event-driven behaviors in web
applications. Modules are compiled into self-contained objects called \emph{reactive machines},
which can be controlled from standard \JavaScript code. It imports core temporal statements from
\Esterel, such as concurrency and preemption, allowing developers to replace complex chains of
callbacks with clear, compositional logic.  By treating external events as synchronous signals,
\textsc{HipHop.js} brings determinism and structured concurrency to an inherently asynchronous
environment.

The key innovation in \textsc{HipHop.js} is the \texttt{async} construct, which bridges the
synchronous machine and the asynchronous world of \JavaScript. An \texttt{async} block is bound to a
signal and behaves as a single, time-consuming statement. When executed, it launches an asynchronous
task and pauses the local synchronous control thread—without blocking other parallel threads—until
the task completes.

\begin{multicols}{2}[The \texttt{AUTHENTICATE} module illustrates the use of \texttt{async} in
    \textsc{HipHop.js}.]
  \begin{lstlisting}[basicstyle=\ttfamily\footnotesize\color{BlackText}]
  hiphop module AUTHENTICATE {
    in NAME, PASSW;
    out STATE, RESULT;
    emit STATE("start");
    async RESULT {
      let req = authenticate(
        NAME.nowval, PASSW.nowval);
      req.on('result',
        res => { this.notify(res) });
      req.end();
    }
    emit STATE("result");
  }
\end{lstlisting}
  \newcolumn
  \begin{table}[H]
    \centering
    \begin{tabular}{l|cccc}
      \texttt{NAME}   & \texttt{"user"}     & $\bot$ & $\bot$            & \dots \\
      \texttt{PASSW}  & \texttt{"password"} & $\bot$ & $\bot$            & \dots \\
      \texttt{RESULT} & $\bot$              & $\bot$ & \texttt{true}     & \dots \\
      \texttt{STATE}  & \texttt{"start"}    & $\bot$ & \texttt{"result"} & \dots \\
    \end{tabular}
    \caption{An execution trace of the \texttt{AUTHENTICATE} module.}
    \label{tab:sota:synchronous:hiphop:exec}
  \end{table}
\end{multicols}

The execution flow, captured by the trace in \Cref{tab:sota:synchronous:hiphop:exec}, unfolds across
several instants. In the first instant, upon receiving \texttt{NAME} and \texttt{PASSW}, the module
emits \texttt{STATE} with the value \texttt{"start"}. It then executes the \texttt{async RESULT} statement,
which launches the external \texttt{authenticate} function and immediately pauses this synchronous
thread. While the network request is in flight, the local thread remains suspended. Finally, in a
future instant, the request completes and its callback invokes \texttt{this.notify}. This triggers a
new reaction where two things happen simultaneously: the \texttt{async} statement emits the
\texttt{RESULT} signal and terminates, allowing the final statement, \texttt{emit STATE("result")},
to execute. Thus, in this final instant, \texttt{RESULT} is present and \texttt{STATE} is updated to
\texttt{"result"}.

\paragraph*{Comparison with GRust}
\GRust shares with \Esterel the goal of providing a formal model for deterministic, reactive
control, but it follows a fundamentally different philosophy. The \GRSync sub-language of \GRust is
declarative and dataflow-oriented, similar to \Lustre, whereas \Esterel is imperative. An engineer
using \GRust describes \emph{what} the relationships between dataflows are, not the sequence of
steps required to compute them. This design aligns \GRust more closely with the block-diagram models
familiar to control engineers. Furthermore, the execution model of \GRust is based purely on the
propagation of signal changes and event occurrences (\emph{push-based}). A significant consequence
is that \GRust cannot detect the absence of an event; computations are only triggered by the arrival
of new data, and without a global clock tick to react against, there is no moment to check if a
signal was absent.

The approach to asynchrony also differs fundamentally. \GRust confines asynchrony to the \GRFRP
layer, which orchestrates asynchronous \emph{dataflows}---that is, streams of values that arrive at
different and irregular intervals. This contrasts with the \texttt{async} construct in
\textsc{HipHop.js}, which serves to integrate a single, time-consuming asynchronous operation into a
synchronous execution thread. The \textsc{HipHop.js} approach is conceptually similar to the futures
in \Lustre~\cite{DBLP:conf/emsoft/CohenGP12}. \GRust adopts a GALS architecture, where the
imperative control of \Esterel is replaced by the declarative composition of dataflow components.

\subsection{Prelude: From Synchronous Models to Multi-Periodic Architectures}
\label{sec:sota:synchronous:prelude}

\Prelude is an architectural language and compiler that builds upon synchronous dataflow languages
like \Lustre, with a specific focus on generating multi-task, hard real-time
systems~\cite{DBLP:journals/deds/PagettiFBCL11}. It directly addresses the inefficiency of the
single-threaded approach common to synchronous compilers, as illustrated by the conceptual \Lustre
code in the example below.

\begin{multicols}{2}
  \begin{lstlisting}[basicstyle=\ttfamily\footnotesize\color{BlackText}]
  node MULTIRATE(I: int)
        returns (O: int)
  var CLK: bool; V: int when CLK;
  let
    O = FAST(I, current(V));
    CLK = true -> not pre(CLK);
    V = SLOW((0 fby O) when CLK);
  tel
\end{lstlisting}
  \newcolumn
  \begin{table}[H]
    \centering
    \begin{tabular}{l|ccccc}
      \texttt{I}   & $i_0$         & $i_1$          & $i_2$         & $i_3$          & \dots \\
      \texttt{CLK} & \texttt{true} & \texttt{false} & \texttt{true} & \texttt{false} & \dots \\
      \texttt{V}   & $v_0$         & $\bot$         & $v_1$         & $\bot$         & \dots \\
      \texttt{O}   & $o_0$         & $o_1$          & $o_2$         & $o_3$          & \dots \\
    \end{tabular}
    \caption{An execution trace of the multi-rate node.}
    \label{tab:sota:synchronous:lustre:gap}
  \end{table}
\end{multicols}

In this model, the \texttt{SLOW} node representes a computation that takes time, and that is only
activated on the clock \texttt{CLK}, which is half the rate of the base clock. A standard compiler
would generate a single \texttt{step} function that must run at the fastest rate. The overall
reaction time of the system would be constrained by the \wcet of the \texttt{SLOW} step function,
forcing the faster \texttt{FAST} node to idle while the slower logic is checked.

\begin{multicols}{2}[\Prelude solves this by providing a language to describe the real-time
    architecture, allowing the compiler to generate a partitioned, multi-task system. The
    following code provides the same system described in \Prelude.]
  \begin{lstlisting}[basicstyle=\ttfamily\footnotesize\color{BlackText}]
  imported node FAST (I, J: int) 
             returns (O: int) wcet 2;
  imported node SLOW (I: int) 
             returns (O: int) wcet 9;

  node MULTIRATE(I: int rate(10, 0))
        returns (O: int)
  var V: int;
  let
    O = FAST(I, current(V) *^ 2);
    V = SLOW((0 fby O) /^ 2);
  tel
\end{lstlisting}
  \newcolumn
  \begin{table}[H]
    \centering
    \begin{tabular}{l|ccccc}
      date       & 0     & 10    & 20    & 30    & \dots \\
      \texttt{I} & $i_0$ & $i_1$ & $i_2$ & $i_3$ & \dots \\
      \texttt{V} & $v_0$ &       & $v_1$ &       & \dots \\
      \texttt{O} & $o_0$ & $o_1$ & $o_2$ & $o_3$ & \dots \\
    \end{tabular}
    \caption{A real-time trace of the \texttt{MULTIRATE} node.}
    \label{tab:sota:synchronous:prelude:exec}
  \end{table}
\end{multicols}

The \Prelude code enriches the synchronous model with real-time properties. The \texttt{wcet} and
\texttt{rate} annotations provide timing information. The key elements are the rate-change
operators: \verb|/^ 2| is a \emph{downsampling} operator that creates a clock twice as slow, while
\verb|*^ 2| is an \emph{upsampling} operator that holds the last value of a slow signal for use by a
faster component.

This enriched specification enables the \Prelude compiler to perform an architectural
transformation. It partitions the node into two clock-coherent tasks: a fast task for \texttt{FAST}
and a slow task for \texttt{SLOW}. It generates a real-time configuration: \texttt{fast\_task} with
a 10ms period and a 2ms WCET, and \texttt{slow\_task} with a 20ms period and a 9ms WCET. It
synthesizes the communication buffers needed to pass data between the tasks, whose sizes are
statically proven to be bounded. This process, formally proven correct, yields a verifiable
multi-task architecture from a high-level model, bridging the gap between logical specification and
real-time implementation.

\paragraph*{Comparison with GRust}
The most significant difference lies in their interpretation of the GALS model. \Prelude implements
a \emph{multi-periodic} GALS, where the asynchronous communication is structured and predictable,
designed to be mapped onto the fixed, periodic tasks of an RTOS. Its goal is to guarantee hard
real-time schedulability.

\GRust, in contrast, implements an \emph{event-driven} GALS. The asynchronous communication at the
\GRFRP level is truly sporadic, designed to react efficiently to an unpredictable environment. Its
target is not a conventional RTOS but a modern asynchronous runtime like \texttt{tokio}.
Consequently, the timing constraints in \GRust (\eg \texttt{debounce}, \texttt{refresh}) are not for
hard real-time scheduling but for managing event load and ensuring liveness. \Prelude brings the
synchronous paradigm to multi-periodic real-time platforms, whereas \GRust embeds synchronous
components within a fully event-driven architecture.

\subsection{Synthesis: GRust in the Synchronous Landscape}
\label{sec:sota:synchronous:synthesis}

The survey of synchronous languages reveals a spectrum of design philosophies, from the declarative
dataflow of \Lustre to the imperative control of \Esterel. \GRust carves out a distinct position by
selectively combining principles from this landscape and adapting them to the specific challenges of
event-driven, high-performance systems. The main distinctions can be summarized along three primary
axes: the programming model, the clocking and execution model, and the target architecture.

First, \GRust aligns firmly with the declarative, dataflow paradigm of \Lustre, not the imperative
style of \Esterel. A \GRust engineer describes relationships between dataflows, which is a natural
fit for the block-diagram models used in control engineering. It deliberately omits the complex
intra-component concurrency and preemption constructs of \Esterel, simplifying the component-level
logic.

Second is its clocking and execution model. Where \Lustre and \Signal offer powerful but complex
clock calculi to manage multi-rate systems, \GRust makes a radical simplification: the \GRSync
sub-language has no user-defined clocks. Instead, it provides only high-level constructs that abstract
away clock manipulation, much like the \texttt{automaton} or \texttt{switch} features in modern
\Lustre. This is coupled with a purely \emph{push-based} execution model, where computations in
\GRFRP are triggered only by the arrival of new data. This design directly addresses the goal of
avoiding redundant computations, but it has a profound consequence: a \GRust component cannot
detect the absence of a signal, as there is no global clock tick to react against.

Finally, these choices are dictated by its target architecture. Languages like \Lustre and \Esterel
target safety-critical, hard real-time environments; \Prelude refines this by mapping synchronous
models to multi-periodic RTOS tasks. \GRust, in contrast, is designed for the event-driven GALS
architecture of modern software-defined vehicles. Its goal is not hard real-time schedulability but
efficient resource management in a high-throughput asynchronous runtime like \texttt{tokio}. This
focus explains why its approach to asynchrony---orchestrating sporadic dataflows at the service
level---differs from the integration of single asynchronous operations found in \textsc{HipHop.js}.

In summary, \GRust can be seen as an adaptation of synchronous principles. It trades the expressive
power of complex clock calculi and imperative control for a simpler, event-driven model that is
formally defined, safe by construction, and tailored for integration into modern, asynchronous
\Rust-based systems.

\section{Languages and Frameworks for Distributed Systems}
\label{sec:sota:distributed}

While synchronous languages provide a powerful model for logically centralized systems, modern
automotive architectures are distributed. This section examines frameworks designed explicitly for
coordinating distributed components, from time-triggered periodic models to best-effort event-based
systems.

\subsection{Giotto: A Time-Triggered Methodology}
\label{sec:sota:distributed:giotto}

\textsc{Giotto} is a programming methodology and language for hard real-time embedded systems that
enforces a strict separation between the functional behavior of tasks and their real-time execution%
~\cite{DBLP:journals/pieee/HenzingerHK03}. A \textsc{Giotto} program is structured around three main
units: tasks, drivers, and modes. The basic functional unit is the \emph{task}, which periodically
executes external code. Tasks communicate with each other, as well as with sensors and actuators,
through \emph{drivers}, which are pieces of code assumed to execute essentially instantaneously,
satisfying the synchrony hypothesis. In this way, the \textsc{Giotto} abstraction integrates
scheduled, time-consuming computations (tasks) with synchronous, instantaneous communication
(drivers). Finally, a set of concurrently executing tasks makes up a \emph{mode}; the system can
switch between modes to change its behavior.

The central concept enabling this is \emph{Logical Execution Time}
(LET)~\cite{DBLP:conf/birthday/KirschS12}. A task reads its inputs precisely at its release time and
produces its outputs precisely at its termination time, regardless of when the underlying
computation actually completes, provided it finishes before its deadline. This decoupling of logical
from physical time is the source of both \textsc{Giotto}'s primary strength and its main drawback.
On one hand, it makes the system's timing behavior perfectly deterministic and predictable: whether
a task runs on a fast or a slow processor, its outputs appear at the same logical moments. On the
other hand, it introduces inherent latency. Because a task's output is not visible until the end of
its logical execution window, the system cannot react faster than this pre-defined duration. This
model, while safe and predictable, can be restrictive and imply high end-to-end latencies from input
acquisition to the corresponding output production.

A \textsc{Giotto} programmer specifies only the logical timing and functional dependencies, it does
not dictate how or when tasks are scheduled. Instead, this responsibility falls to the compiler. It
translates the program---composed of tasks, modes, and their periodic schedules---into a
time-triggered instruction set, performing the necessary schedulability analysis to generate a
static schedule that guarantees all timing constraints can be met on the target platform.

\paragraph*{Comparison with GRust}
The primary distinction between \textsc{Giotto} and \GRust lies in the philosophy of
\emph{time-triggered} versus \emph{event-triggered} execution. \textsc{Giotto} is
fundamentally time-triggered: computations happen at pre-determined points in time, dictated
by a static schedule, regardless of whether inputs have changed. \GRust, conversely, is
fundamentally event-triggered, with computations happening only in response to the arrival of
new data. This leads to a key trade-off in the context of the \sdv: \textsc{Giotto} provides
predictable timing for hard real-time control at the potential cost of inefficiency (running
computations on unchanged data), while \GRust prioritizes efficiency by avoiding redundant
computations, giving up the hard real-time guarantees that come from a static,
time-triggered schedule.

\subsection{Lingua Franca: A Reactor-Oriented Coordination Language}
\label{sec:sota:distributed:lf}

The design of \textsc{Lingua Franca} (LF) is heavily influenced by actor-oriented frameworks like
\emph{Ptolemy~II}~\cite{DBLP:journals/pieee/EkerJLLLLNSX03}. Building on these concepts, LF is a
polyglot coordination language based on \emph{reactors}---components with input and output ports and
\emph{reactions} (code blocks) that execute in response to events carrying logical timestamps, or
\emph{tags}~\cite{DBLP:journals/tcad/LeeS98}. A dedicated runtime enforces a deterministic execution
order based on these tags and a dependency analysis of the reactions, guaranteeing that the behavior
of the system is repeatable~\cite{DBLP:journals/taco/MenardLBCDDFLSTKCL23}. LF preserves determinism
by default, admitting nondeterminism only when explicitly introduced by the programmer, and can
target multiple languages, including \Ci, \Cpp, \textsc{Python}, \textsc{TypeScript}, and \Rust.

The \texttt{Sample} reactor below, which targets \Rust, illustrates these concepts.

\begin{lstlisting}[basicstyle=\ttfamily\footnotesize\color{BlackText}]
target Rust
reactor Sample(offset: time = 0, period: time) {
  input i: i32
  output o: i32
  state s: i32 = 0
  timer t(offset, period)
  reaction(t) -> o {= ctx.set(o, ctx.get(s)); =}
  reaction(i) {= self.s = ctx.get(i).unwrap(); =}
}
\end{lstlisting}

This reactor defines two event-triggered computations, with native \Rust code specified inside
\texttt{\{= ... =\}} blocks. The first reaction is triggered by an event on the input port
\texttt{i} and updates the internal \texttt{state s} of the reactor. The second is triggered by a
periodic \texttt{timer t} and sets the output port \texttt{o} on the current state. This setup
creates a classic sampler: the state is updated whenever new input arrives, but the output is only
produced periodically. The LF runtime guarantees that if a timer event and an input event occur at
the same logical time, the input reaction will execute first due to the implicit data dependency
through the state variable \texttt{s}.

Where a language like \Prelude partitions a single synchronous program into a fixed set of
\emph{periodic} tasks, LF provides a more general coordination layer for any set of event-driven
components. It ensures deterministic execution even with sporadic events, separating the
coordination logic from the computational code within the reactions.

\paragraph*{Comparison with GRust}
Both LF and the \GRFRP layer of \GRust can be seen as coordination languages for event-driven
components, but their core guarantees and programming models differ. This is most evident in
how they define behavior. In LF, the programmer explicitly defines the chain of causality by
specifying reactions and their triggers---in effect, declaring the event propagation path. In
\GRust, the programmer uses a declarative, equation-based style to define the relationships between
dataflows, and it is the compiler's responsibility to derive the propagation strategy.

This distinction leads to their different core guarantees. The primary contribution of LF is
\emph{deterministic concurrency}: its centralized, time-aware runtime guarantees a globally
consistent and repeatable order of reactions. \GRust, by targeting an asynchronous runtime like
\texttt{tokio}, operates in a best-effort asynchronous world and does not guarantee a deterministic
execution order \emph{between} its services. Instead, \GRust focuses on providing formal dataflow
correctness, determinism, and memory safety \emph{within} each service.

\subsection{Robot Operating System (ROS)}
\label{sec:sota:distributed:ros}

The \emph{Robot Operating System} (ROS) is a widely adopted open-source framework and middleware for
building complex, distributed robotic systems~\cite{ros}. Its architecture is based on a
decentralized, publish-subscribe model. A system is composed of independent processes called
\emph{nodes}, which communicate by passing messages over named channels called \emph{topics}.

However, the default communication model in ROS is fundamentally asynchronous and non-deterministic.
There is no central clock or coordinator, and message delivery is best-effort, with no built-in
guarantees on latency or ordering. While this flexibility is excellent for rapid prototyping, it
makes formal reasoning about system-wide properties, especially timing, extremely difficult.
Correctness relies on careful node-level design and extensive integration testing rather than on
formal guarantees from the framework itself.

The Python code below, implementing a simple robot controller, highlights a key consequence of
this model.

\begin{lstlisting}[basicstyle=\ttfamily\footnotesize\color{BlackText}]
from rospy import Publisher, Subscriber, init_node, spin
from std_msgs.msg import Float64
import threading

class RobotController:
  def __init__(self):
    self.lock = threading.Lock()
    self.x = 0.0
    self.goal = 0.0
    self.cmd_pub = Publisher('/robot/cmd', Float64, queue_size=10)
    self.x_sub = Subscriber('/robot/x', Float64, self.x_callback)
    self.goal_sub = Subscriber('/robot/goal', Float64, self.goal_callback)

  def x_callback(self, msg: Float64):
    with self.lock:
      self.x = msg.data
      self._move()

  def goal_callback(self, msg: Float64):
    with self.lock:
      self.goal = msg.data
      self._move()

  def _move(self):
    error = self.goal - self.x
    if abs(error) >= 0.01:
      self.cmd_pub.publish(error)
    else:
      self.cmd_pub.publish(0.0)

if __name__ == '__main__':
  init_node('robot_controller')
  controller = RobotController()
  spin()
\end{lstlisting}

Functionally, the node implements a simple proportional controller to move a robot to a goal
position. Architecturally, it reveals a critical challenge. The callbacks \texttt{x\_callback} and
\texttt{goal\_callback} can be executed concurrently by the ROS runtime. Without protection, this
creates a race condition: one callback could be preempted by the other after reading a shared state
variable but before acting on it, leading to a command being issued based on inconsistent,
partially-updated state. To prevent this non-deterministic behavior, the programmer must manually
introduce a \texttt{threading.Lock}. This explicit, manual management of concurrency is notoriously
difficult and error-prone, yet it is essential for achieving safety and determinism in a standard
ROS node.

\paragraph*{Comparison with GRust}
At an architectural level, the \GRust model of services communicating \via imports and exports is
conceptually similar to ROS nodes communicating over topics. Both frameworks embrace an
event-driven, asynchronous philosophy for system-level coordination. The key difference lies in the
guarantees provided for the components themselves. A ROS node is typically an unstructured program
where the burden of ensuring correctness and concurrency safety falls entirely on the developer. In
contrast, a \GRust service is constructed from formally defined, synchronous \GRSync components.
\GRust thus brings the rigor of synchronous programming and the safety of the \Rust language to the
implementation of individual components within a larger, ROS-like asynchronous architecture,
abstracting away the need for manual locking and imperative propagation logic.

\subsection{Synthesis: GRust in the Distributed Systems Landscape}
\label{sec:sota:distributed:synthesis}

The frameworks surveyed for distributed systems reveal a fundamental design spectrum, with the
trade-off between determinism and efficiency at its core. At one end lies the \emph{time-triggered}
model of \textsc{Giotto}, which provides the strongest guarantees of predictability and temporal
determinism through a static, pre-verified schedule. This safety, however, comes at the cost of
potential inefficiency, as computations are performed at fixed intervals, regardless of whether new
data is available.

At the other end of the spectrum is the \emph{best-effort, event-triggered} model of ROS. It offers
maximum flexibility and responsiveness, making it ideal for rapid prototyping and loosely-coupled
systems. However, it provides no built-in guarantees of determinism, latency, or ordering, placing
the entire burden of ensuring safety and managing concurrency on the developer, often through
complex and error-prone manual locking.

\textsc{Lingua Franca} carves out a third, distinct position. It is event-triggered but not best-effort. By
embedding logical timestamps into events and using a centralized, time-aware runtime, it achieves
\emph{deterministic concurrency}, combining the responsiveness of an event-driven model with the
repeatability of a time-triggered one.

\GRust offers another distinct point on this spectrum, one that is pragmatically tailored to the
modern automotive context. Architecturally, it embraces a similar event-driven, best-effort model to
ROS at the system level, acknowledging that the underlying hardware and network of an \sdv are
inherently asynchronous. However, its philosophy diverges at the component level. Instead of
enforcing global determinism like \textsc{Lingua Franca}, or leaving it entirely to the programmer like ROS,
\GRust focuses on guaranteeing \emph{local determinism} at the service level.

\section{Functional Reactive Programming}
\label{sec:sota:frp}

Functional Reactive Programming (\frp) offers a declarative approach to programming reactive systems
by modeling behavior as functions over time-varying values. Originating in the functional
programming community, it has evolved into a broad family of languages and libraries, each with
different trade-offs between expressivity, performance, and safety guarantees.

\subsection{Classical and Arrowized FRP}
\label{sec:sota:frp:classical}

Classical \frp, as introduced in \textsc{Fran}~\cite{DBLP:conf/icfp/ElliottH97}, is defined by two
core abstractions: \emph{Behaviors}, representing continuous time-varying values, and \emph{Events},
representing discrete occurrences in time. Programs are constructed by composing pure functions over
these signals, enabling a highly declarative and compositional style.

A canonical example is the modeling of a bouncing ball, whose state changes discretely (at each
bounce) while its motion is described by continuous physics between bounces.

\begin{lstlisting}[basicstyle=\ttfamily\footnotesize\color{BlackText}]
bouncing y0 v0 t0 = start t0 (y0, v0)
  where
    start t (y, v) = motion `untilB' bounce +=> start
      where
        motion = pairB pos vel
        pos    = lift0 y + integral vel
        vel    = lift0 v + integral (-9.81)
        bounce = ((predicate (pos <=* 0) t) `snapshot' motion)
                   ==> \(_, (yB, vB)) -> (yB, -0.8*vB)
\end{lstlisting}

This code defines a recursively specified behavior. Let us break down the \frp concepts it uses.
The position (\texttt{pos}) and velocity (\texttt{vel}) are \emph{Behaviors}. They are defined using
continuous-time concepts like \texttt{integral}, which describes the physics of motion under
constant gravity. The \texttt{motion} is a behavior that pairs them together. An \emph{Event} is
created using a \texttt{predicate} that monitors a boolean \emph{Behavior}. Here, an event is
triggered whenever the ball hits the floor (\texttt{pos <=* 0}). When the collision event from the
\texttt{predicate} occurs, \texttt{snapshot} captures the \emph{current value} of the
\texttt{motion} behavior at that exact moment. The \texttt{untilB} operator switches from one
behavior to another when an event occurs. The program follows the continuous \texttt{motion}
behavior \emph{until} the \texttt{bounce} event. The \texttt{+=>} operator then defines the next
behavior. It takes the value from the \texttt{bounce} event (which contains the new, damped
velocity) and uses it to recursively call \texttt{start}, defining the next phase of continuous
motion.

However, this expressive denotational model, based on lazy evaluation, suffered from practical
implementation challenges. The purity of the model made it difficult for the runtime to garbage
collect the history of signals, leading to space and time leaks where resource consumption would
grow without bound~\cite{DBLP:conf/icfp/Krishnaswami13}.

To address these limitations, \emph{Arrowized \frp} (AFRP), as implemented in
\Yampa~\cite{DBLP:conf/haskell/CourtneyNP03}, introduced a more structured approach. AFRP
makes the static dataflow structure of the program explicit. Programs are built from \emph{signal
  functions} (type \texttt{SF a b}), which are stateful transformations from an input signal of
type \texttt{a} to an output signal of type \texttt{b}. The \texttt{ballGUI} example below
demonstrates how these are composed using arrow syntax.

\newpage

\begin{lstlisting}[basicstyle=\ttfamily\footnotesize\color{BlackText}]
mouseSF  :: SF GUIInput Point
moveBall :: Point -> Picture
ballGUI  :: SF GUIInput Picture
ballGUI = proc gin -> do
               mouse <- mouseSF         -< gin
               bpic  <- liftSF moveBall -< mouse
               returnA -< bpic
\end{lstlisting}

Here, the program is a signal function that transforms GUI inputs into a picture. The \texttt{proc}
block enables arrow notation, where signal functions are wired together. The \texttt{mouseSF} signal
function takes the raw GUI input signal \texttt{gin} and produces a signal of mouse coordinates
(\texttt{Point}). The \texttt{-<} operator is the application that wires the input signal to the
signal function. The \texttt{liftSF} combinator lifts a pure function (\texttt{moveBall}) into a
stateless signal function. This new signal function is then applied to the \texttt{mouse} signal to
produce the final picture signal, \texttt{bpic}. \texttt{returnA} specifies the final output signal
of the composed \texttt{ballGUI} signal function. This arrow syntax makes the dataflow explicit,
allowing the compiler to manage state within each signal function and guarantee bounded resource
usage, at the cost of some dynamic flexibility.

\paragraph*{Comparison with GRust}
The design of \GRust is directly informed by this history. It avoids the problems of classical \frp
by imposing a strict separation of concerns: the \GRFRP sub-language is a \emph{compositional} layer for
wiring together dataflow operators, while new operators must be defined in the \GRSync sub-language. The
\GRSync components are conceptually similar to the signal functions of \Yampa, providing the
same guarantee of bounded resource usage.

However, \GRust goes further to address a subtle but fundamental tension in the handling of feedback
loops in \Yampa with the \texttt{loop} operator. In a purely push-based (on-change)
execution, it is ambiguous how a stateful feedback loop should be evaluated. If it re-evaluates on
every input change, the result may be semantically incorrect, as the output of the stateful node
depends on its internal state, not just the current input. If it evaluates periodically, it loses
the efficiency of the push-based model. If it computes a fix-point, it may not terminate.

\GRust addresses this challenge directly with what we call a \emph{reactive type system}
(\Cref{sec:semantics:grsync:react_typing,sec:semantics:grfrp:react_typing}). This
system statically ensures that for any well-typed service, a purely push-based execution (reacting
only to input changes) is semantically equivalent to a fine-grained synchronous periodic execution
(\Cref{thm:semantics:grfrp:theorems:oversampling}). This provides a formal guarantee that the
efficient, event-driven execution strategy is also correct, resolving the ambiguity present in other
push-based systems.

\subsection{Real-Time and Event-Driven FRP}
\label{sec:sota:frp:rt}

To make \frp suitable for resource-constrained and real-time environments, researchers have proposed
variants that impose restrictions to gain predictability. Early approaches focused on imposing
semantic restrictions on the classical model. \emph{Real-Time \frp}
(RT-\frp)~\cite{DBLP:conf/icfp/WanTH01} ensures bounded memory and time by discretizing signals and
restricting recursive definitions, making the signal history finite and manageable.
\emph{Event-Driven \frp} (E-\frp)~\cite{DBLP:conf/padl/WanTH02} addresses causality concerns by
making all updates explicitly triggered by discrete events, ensuring a clear and deterministic
signal propagation path.

To make the concepts of RT-\frp and E-\frp more concrete, we can examine how each would model a
simple reactive computation. The syntax used here is hypothetical but reflects the core principles
of each approach.

\paragraph*{RT-FRP: Bounded State through Discretization}
The goal of RT-\frp is to guarantee bounded time and memory. It achieves this by treating
signals as discrete sequences (like in synchronous languages) and by restricting recursion to
occur only across time steps. Consider a simple event counter.

\begin{lstlisting}[basicstyle=\ttfamily\footnotesize\color{BlackText}]
update_signal = let snapshot inc <- IncButton in
  let snapshot dec <- DecButton in
    (ext if is_some(inc) then 1 else if is_some(dec) then -1 else 0)

count_signal = let snapshot update <- update_signal in
  let snapshot count <- delay 0 (ext count + update) in 
    (ext count)
\end{lstlisting}

This model operates in discrete, synchronous steps. Let's break down the syntax. The
\texttt{let snapshot x <- s} takes the current value of the signal \texttt{s} and stores it as
\texttt{x}. The \texttt{update\_signal} is thus a stream whose value is \texttt{1}, \texttt{-1}, or
\texttt{0} at each instant. The definition of \texttt{count\_signal} uses the expression
\texttt{delay 0 (...)}, a direct analog to the \texttt{pre} operator in \Lustre. It initializes the
signal with \texttt{0} and, in every subsequent step, provides the value of the signal from the
previous step.

By enforcing that recursion must be guarded by a \texttt{delay}, RT-\frp ensures that the
computation for each step is finite and non-circular. The system only needs to store the state from
the single previous step to compute the next, guaranteeing bounded memory and solving the space leak
problem of classical \frp.

\paragraph*{E-FRP: Guaranteed Causality through Event Triggers}
The goal of E-\frp is to ensure causality: an output at time $t$ can only depend on inputs that
occurred at or before $t$. It achieves this by ensuring all state changes are explicitly triggered
by discrete events.

\begin{lstlisting}[basicstyle=\ttfamily\footnotesize\color{BlackText}]
count = init x = 0 in
    { IncButton => x + 1 }
    { DecButton => x - 1 }

color = if count < 0 then (0, 0, 255) 
        else if count >= 100 then (255, 0, 0) else (0, 255, 0)
\end{lstlisting}

This model is not based on continuous-time behaviors or synchronous steps, but on atomic state
updates. The program defines a piece of state, \texttt{count}, initialized to \texttt{0}. It then
declares a set of \emph{event handlers}. When an \texttt{IncButton} event arrives, the state of
\texttt{count} is updated by applying the pure function \texttt{x + 1} to the current state. The
same logic applies to \texttt{DecButton}. The \texttt{color} value is a pure, stateless function
derived from the current value of the \texttt{count} state.

This event-handler style makes causality explicit. The state of the system evolves in a series of
discrete, atomic transitions, each triggered by a specific input event. There is no ambiguity about
when a computation happens, which was a key challenge in classical \frp.

\paragraph*{Comparison with GRust}
\GRust can be seen as a synthesis of the principles found in both RT-FRP and E-FRP, but with a
strict architectural separation.

On one hand, the \GRSync sub-language is conceptually identical to RT-FRP. It is a discrete, synchronous
dataflow model where recursion is guarded by a delay operator (\texttt{last}) to guarantee bounded
time and memory. A \GRSync component is, in essence, a formally-defined RT-FRP signal function,
providing the same strong guarantees on resource usage and determinism.

On the other hand, the \GRFRP service layer embodies the philosophy of E-FRP. It is a purely
event-driven system where computations are triggered by the arrival of asynchronous events, ensuring
a clear causal link from input to computation.

The key innovation of \GRust lies not in choosing one model over the other, but in its explicit
integration of both within a Globally Asynchronous, Locally Synchronous (GALS) architecture. Where
RT-FRP and E-FRP are distinct language-level solutions to the problems of classical FRP, \GRust uses
their principles as complementary architectural layers. The result is a system that provides the
formal, analyzable safety of synchronous programming for its core computations (\GRSync) while
retaining the efficiency and responsiveness of a fully event-driven system at the service level
(\GRFRP).

\subsection{Modal FRP for Safe Reactive Programming}
\label{sec:sota:frp:modal}

While the restrictions of RT-\frp and E-\frp are effective, a more recent and principled approach is
\emph{Modal \frp}, which uses advanced type systems to enforce resource safety by
construction~\cite{DBLP:conf/icfp/Krishnaswami13,DBLP:conf/popl/KrishnaswamiBH12,DBLP:conf/sfp/DischHB25}.
\textsc{Rattus}~\cite{DBLP:journals/pacmpl/BahrGM19} is an example of this approach. It employs a
modal type theory to distinguish between computations that can be executed in the current time step
(\emph{now}) and those that must be delayed until the \emph{next} time step. This distinction,
enforced by the type checker, statically prevents common \frp bugs like space leaks and causality
violations.

In a synchronous modal \frp like \textsc{Rattus}, the type system provides strong guarantees. The
core idea is to use a modal type, written \texttt{O a}, to represent a computation that will produce
a value of type \texttt{a} in the next time step. The type checker ensures that any recursive
definition on a stream must pass through this \texttt{O} modality (\eg \via a \texttt{delay}
operator). This rule makes it impossible to define a stream that depends on its own value at the
same instant, thus guaranteeing productivity and causality.

However, the synchronous model, which relies on a single global clock, is inefficient for many
applications. For example, in a GUI, a text field needs to react to keyboard input, while a button
reacts to mouse clicks. A global clock would force both components to execute a computation step for
every single input, wasting resources. To address this, the extension
\textsc{Async Rattus}~\cite{DBLP:journals/pacmpl/BahrM23} adapts the modal model with dynamic, local
clocks, enabling efficient asynchronous computation.

\begin{lstlisting}[basicstyle=\ttfamily\footnotesize\color{BlackText}]
map :: Box (a -> b) -> Sig a -> Sig b
map f (x ::: xs) = unbox f x ::: delay (map f (adv xs))

const :: a -> Sig a
const x = x ::: never

switch :: Sig a -> O (Sig a) -> Sig a
switch (x ::: xs) ysd = x ::: delay (case select xs ysd of
                                       Fst  xs' ysd' -> switch xs' ysd'
                                       Snd  _   ys   -> ys
                                       Both _   ys   -> ys)
\end{lstlisting}

In \textsc{Async Rattus}, the semantics of the \texttt{O} modality changes: it represents a delay
relative to a \emph{local} clock, not a global one. The above code, taken from the signal library of
\textsc{Async Rattus}~\cite{async_ratt}, illustrates key concepts. The recursive call in
\texttt{map} is guarded by \texttt{delay}, which is required by the type system to ensure
productivity. The function itself is passed in a \texttt{Box}, representing a time-invariant,
reusable computation. The \texttt{const} signal uses \texttt{never}, a computation whose local clock
will never tick, creating a truly static value. The \texttt{switch} function allows a signal to
behave like \texttt{xs} initially, but dynamically change its behavior to that of a new signal
\texttt{ys} as soon as the local clock of \texttt{ys} ticks.

\paragraph*{Comparison with GRust}

\GRust is philosophically aligned with the goal of Modal \frp of achieving safety by construction,
but it draws its solutions from a different research history: the established principles of
synchronous dataflow languages.

Causality and productivity are not enforced through a modal type system. Instead, they are
guaranteed by the simpler dataflow structure of \GRSync components, where the \texttt{last} operator
provides a straightforward, non-modal form of delay that is sufficient for static causality
analysis. Similarly, memory safety is not derived from a custom type theory. \GRust leverages two
powerful, orthogonal mechanisms: synchronous compilation ensures that each component has a
statically bounded memory footprint, while the underlying \Rust compiler provides compile-time
guarantees against memory errors \via its ownership and borrow-checking system.

This design allows \GRust to offer strong safety guarantees without requiring its users to
understand modal logic, making it more accessible to automotive engineers.

\subsection{Practical FRP-Inspired Ecosystems}
\label{sec:sota:frp:practical}

The core ideas of \frp have been widely adopted through a generation of libraries that, while
inspired by \frp, often deviate significantly from its classical formulation. Systems like the
\ReactiveX (Rx) family~\cite{reactivex} focus on discrete event streams and adopt a predominantly
\emph{push-based} or hybrid execution model. This makes them highly responsive and a natural fit for
handling user input and network events in modern applications. Their primary strength lies in
providing a rich set of compositional operators (\eg \texttt{map}, \texttt{filter}, \texttt{merge},
\texttt{scan}) that allow developers to manage complex asynchronous dataflows declaratively.

The \ReactiveX (RxJS) code below illustrates this style. It defines a declarative pipeline that
transforms a stream of events based on timing.

\begin{lstlisting}[basicstyle=\ttfamily\footnotesize\color{BlackText}]
var times = [ { val: 0, t: 100 }, { val: 1, t: 600 },
              { val: 2, t: 400 }, { val: 3, t: 700 } ];
var source = Rx.Observable.from(times)
  .flatMap((item) => Rx.Observable.of(item.val).delay(item.t));
var double = source.debounce(500).map((x) => 2*x);
\end{lstlisting}

Let's trace the execution of this code. The \texttt{flatMap} operator processes the initial
observable resulting from the array. For each item, it creates a new observable that will emit the
value of the item after a specified delay. All these timers start at roughly the same time. This
results in a \texttt{source} stream that will emit values on the following approximate timeline:
\begin{itemize}
  \item At t = 100\,ms: emits \texttt{0}
  \item At t = 400\,ms: emits \texttt{2}
  \item At t = 600\,ms: emits \texttt{1}
  \item At t = 700\,ms: emits \texttt{3}
\end{itemize}
This \texttt{source} stream is then piped into the \texttt{debounce(500)} operator. This operator
only emits a value if a 500\,ms period of silence passes after that value was received. If a new
value arrives before the 500\,ms is up, the timer is reset and the previous value is discarded. The
value \texttt{0} (at 100\,ms) is discarded because \texttt{2} arrives 300\,ms later. The value
\texttt{2} (at 400\,ms) is discarded because \texttt{1} arrives 200\,ms later. The value \texttt{1}
(at 600\,ms) is discarded because \texttt{3} arrives 100\,ms later. Only the final value,
\texttt{3}, is followed by a quiet period longer than 500\,ms. Therefore, approximately 1200\,ms
from the start (700\,ms + 500\,ms), the \texttt{debounce} operator emits \texttt{3}. This value is
then passed to the \texttt{map} operator, and the final \texttt{double} stream emits the value
\texttt{6}. This example demonstrates how a complex, stateful, and time-dependent behavior can be
expressed concisely by composing a few powerful operators.

While powerful for data transformation, this style can lead to challenges when managing the overall
state of a complex application, as state can become implicitly distributed across many
interconnected streams. The evolution of the \Elm language is particularly instructive in
this context~\cite{DBLP:conf/pldi/CzaplickiC13}. \Elm began as a web-focused \frp language
with a pure signal-graph model, but later transitioned to a more explicit state management
architecture.

Consider a simple counter application with two buttons, \texttt{incButton} and \texttt{decButton}.
In early, \frp-style \Elm, this would be modeled by creating two independent event streams
(one for each button) and merging them into a single stream of \texttt{updates}.

\begin{lstlisting}[basicstyle=\ttfamily\footnotesize\color{BlackText}]
incEvents = Signal.map (\_ -> 1) incButton.clicks
decEvents = Signal.map (\_ -> -1) decButton.clicks
updates = Signal.merge incEvents decEvents
count = Signal.foldp (+) 0 updates
main = Signal.map showCount count
\end{lstlisting}

In this model, the state (\texttt{count}) is not stored in a variable but exists implicitly as the
result of folding all past update events from the merged stream. While elegant, this approach proved
difficult to scale. In large application, the state of the application should be passed through
folding operations.

To solve this, \Elm moved away from the signal graph to \emph{The {Elm} Architecture}
(TEA), a more predictable, centralized state management pattern.

\begin{lstlisting}[basicstyle=\ttfamily\footnotesize\color{BlackText}]
type alias Model = Int
type Msg = Increment | Decrement
init = 0
update msg model =
  case msg of
    Increment -> model + 1
    Decrement -> model - 1
view model =
  div []
    [ button [ onClick Increment ] [ text "+" ],
      div [] [ text (String.fromInt model) ],
      button [ onClick Decrement ] [ text "-" ] ]
main = Browser.sandbox { init = init, update = update, view = view }
\end{lstlisting}

In TEA, the entire application state is held in a single, explicit \texttt{Model}. The only way to
change the state is to dispatch a \texttt{Msg}, which is processed by a pure \texttt{update}
function. The \texttt{view} is also a pure function that renders the current \texttt{Model}. This
shift created a clear, unidirectional dataflow and a single source of truth, making applications far
easier to reason about and debug. This evolution highlights a broader trend: for mainstream
application development, many ecosystems have favored simpler, more predictable state management
models over the full, unconstrained power of classical \frp.

\paragraph*{Comparison with GRust}
\GRust is strongly aligned with this pragmatic family of languages. The \GReact library is
explicitly inspired by the operators in \ReactiveX. The \GRFRP service layer, which composes
operators on event streams, directly mirrors the declarative style of these libraries.

\Elm addressed the challenges of implicit, distributed state by adopting a single,
centralized state model (The \Elm Architecture). \GRust provides a more compositional
solution to the same problem. The crucial difference lies in what constitutes a user-defined
operator. In a system like \ReactiveX, a custom operator is typically an unstructured function in
the host language (\eg \JavaScript), which can implicitly hold state and have arbitrary side
effects.

In \GRust, a user-defined operator must be a formally-defined synchronous \GRSync component. Each
component is effectively an isolated, synchronous state machine, with a well-defined state that
evolves predictably from one discrete instant to the next. This design choice is central to our
contribution: it provides a formal foundation for building \emph{stateful} operators, ensuring that
even complex, state-dependent logic remains analyzable and predictable. In essence, while
\Elm enforces discipline at the application level, \GRust enforces it at the component
level, allowing for safe composition of complex stateful behaviors.

\subsection{Synthesis: GRust in the Functional Reactive Landscape}
\label{sec:sota:frp:synthesis}

The history of \frp is a compelling narrative of trade-offs between expressive power, formal safety,
and practical usability. Classical \frp offered unparalleled declarative elegance but was plagued by
resource leaks. Arrowized and Safe \frp variants solved these leaks, either by restricting the
programming model or by introducing advanced type systems, but often at the cost of increased
complexity. In parallel, practical libraries like \ReactiveX gained massive adoption by prioritizing
a simple, push-based model, while frameworks like \Elm ultimately pivoted away from implicit
signal graphs entirely toward a more predictable, centralized state architecture.

Like Arrowized and Safe \frp, \GRust guarantees bounded memory and time. However, it achieves this
not through a modal type system, but through its GALS architecture. The \GRSync sub-language enforces the
same ``delayed recursion" principle as RT-\frp and synchronous languages, ensuring that any stateful
component is provably resource-bounded. This provides the safety of academic \frp using a simpler,
synchronous mechanism.

The evolution of \Elm highlights the difficulty of managing state in a purely compositional,
signal-based model. The \Elm Architecture (TEA) solves this with a single, monolithic state
object. \GRust offers a more compositional alternative. By requiring all user-defined operators to
be synchronous \GRSync components, it ensures that each piece of stateful logic is an isolated,
well-defined state machine. This allows for the safe composition of complex, stateful behaviors.

Practical libraries like \ReactiveX are efficient but lack a clear formal semantics, especially
regarding stateful feedback loops. \GRust embraces the efficient, push-based execution model but
formalizes it. The reactive type system (\Cref{sec:semantics:grsync:react_typing,sec:semantics:grfrp:react_typing})
guarantees that for any well-typed service, this
event-driven execution is semantically equivalent to a more easily understood periodic execution.
This provides the efficiency of practical \frp with the correctness guarantees of formal synchronous
models.

In essence, \GRust navigates the \frp landscape by combining the compositional, event-driven style
of practical libraries at the service level with the formal rigor and guaranteed safety of
synchronous languages at the component level.

\section{Verification of Synchronous Dataflow Programs}
\label{sec:sota:verification}

The development of critical systems requires not only a sound programming model but also a rigorous
methodology for verification and analysis. This section surveys techniques and tools used to
establish confidence in system correctness.

\subsection{Property Checking with Observers and Contracts}
\label{sec:sota:verification:properties}

A primary goal of formal verification is to prove that a system satisfies key safety properties. In
the context of synchronous languages, this is commonly achieved using two main approaches: the
observer pattern and contract-based verification.

The \emph{observer pattern}~\cite{TR-2018-11} encodes a property as a synchronous node that monitors
the inputs and outputs of the system under verification. It typically consists of two parts: a node
for the assumed hypotheses and a node for the expected properties. The observer outputs a boolean
stream that is true as long as the property holds and becomes false at the first violation. The
verification task is then transformed into proving that this output can never be false, a proof
typically discharged by a model checker such as the BDD-based tool \textsc{Lesar}~\cite{DBLP:phd/hal/Ratel92}.

\begin{lstlisting}[basicstyle=\ttfamily\footnotesize\color{BlackText}]
node counter(sec, bea: bool) returns (ontime, late, early: bool);
var diff: int;
let
  diff   = (0 -> pre(diff)) + (if bea then  1 else 0)
                            + (if sec then -1 else 0);
  early  = diff <= 2;
  late   = diff >= -2;
  ontime = not (early or late);
tel

node assume(sec, bea: bool) returns (hypothesis: bool);
var forbid_sec, forbid_bea: bool;
let
  forbid_sec = if (false -> pre(forbid_sec)) then not bea else sec;
  forbid_bea = if (false -> pre(forbid_bea)) then not sec else bea;
  hypothesis = not(sec and forbid_sec or bea and forbid_bea);
tel

node expect(ontime: bool) returns (ok: bool);
let
  ok = ontime
tel


node verification(sec, bea: bool) returns (hypothesis, ok: bool);
var ontime, late, early: bool;
let
  ontime, late, early = counter(sec, bea);
  hypothesis = assume(sec, bea);
  ok = expect(ontime);
tel
\end{lstlisting}

The example above, from~\cite{TR-2018-11}, monitors the speed of a train by counting beacon
(\texttt{bea}) and second (\texttt{sec}) signals. The property to verify is ``if the train maintains
cruising speed, it remains on time." The \texttt{assume} node formalizes the hypothesis that beacons
and seconds alternate. A model-checker like \textsc{Lesar} can then formally verify that if this
\texttt{hypothesis} is always true, the property \texttt{ok} also remains true.

A complementary approach, particularly suited for SMT-based tools like
\Kind~\cite{DBLP:conf/cav/ChampionMST16}, is \emph{contract-based verification}. Properties are
specified as annotations directly within the code using pre-conditions (\texttt{assume}) and
post-conditions or invariants (\texttt{guarantee}).

\begin{lstlisting}[basicstyle=\ttfamily\footnotesize\color{BlackText}]
node avg(x, y: real) returns (a: real);
(*@contract
  assume x <= y;
  guarantee x <= a and a <= y; *)
let
  a = (x + y) / 2.0;
tel

node max(x: real) returns (m: real);
let
  m = x -> if x > pre x then x else pre m;
tel

node sav(x: real) returns (s: real);
(*@contract
  assume x > 0.0 and x > pre x;
  guarantee s <= max(x); *)
let
  s = avg(x -> pre s, x);
tel
\end{lstlisting}

The code above, from~\cite{DBLP:conf/cav/ChampionMST16}, illustrates this method. The node
\texttt{sav} defines a stateful smoothing average. The contract specifies that if we \texttt{assume}
the input stream \texttt{x} is positive and strictly increasing, we want to \texttt{guarantee} that
the output \texttt{s} will never exceed the maximum value of \texttt{x} seen so far. A tool like
\Kind translates the node and its contract into a set of logical formulas and uses an SMT solver to
prove the property, often employing techniques like k-induction~\cite{DBLP:conf/fmcad/SheeranSS00}
to reason about the infinite execution of the synchronous program.

\paragraph*{Comparison with GRust}
\GReusot, the verification tool for \GRust, implements the contract-based approach. \GReusot adopts
a similar annotation style to allow developers to specify pre- and post-conditions on \GRSync
components. However, its methodology differs in a crucial aspect: the verification target. While
\Kind verifies the abstract synchronous model, \GReusot translates its specifications into
annotations \emph{on the generated code} for \Creusot, a verifier for the \Rust language itself.
This means the formal proof applies directly to a low-level artifact, bridging the gap between the
model and the implementation.

This choice of targeting \Rust provides a second, significant advantage. A verifier for a \Ci-like
language must contend with aliasing, pointer arithmetic, and memory management, requiring extensive
annotations to prove memory safety before even considering functional correctness. \Creusot, in
contrast, leverages the ownership and borrow-checking system of the \Rust compiler. Since the
compiler has already statically guaranteed memory and thread safety, the verifier can focus
exclusively on the logical correctness of the algorithm, greatly simplifying the proof effort
required from the developer.

\subsection{Heptagon/BZR: Contract-Based Synthesis}
\label{sec:sota:synchronous:heptagon}

\Heptagon/BZR is a \Lustre-like language that integrates the concept of \emph{behavioral contracts}
directly into the language, not just for verification, but for the purpose of
\emph{synthesis}~\cite{DBLP:conf/lctrts/DelavalMR10}. A contract specifies assumptions about the
inputs of a component and guarantees about its outputs. Crucially, it can also declare
\emph{controllable} variables, for which the compiler must synthesize control logic to ensure the
contract is never violated.

This approach shifts the developer's burden from manual implementation and post-hoc verification to
the formal specification of requirements. The compiler leverages these contracts to automatically
generate a correct-by-construction controller, enabling a higher degree of automation in the design
process.

\begin{lstlisting}[basicstyle=\ttfamily\footnotesize\color{BlackText}]
node twomodes (v:int) = (y:int)
  contract
    assume  (v <= 1)  & (v >= 0)
    enforce (o <= 10) & (o >= 0)
    with    (c:bool)
var last x : int = 0;
let
  y = x;
  automaton
    state Up
      do x = last x + v
      until c then Down
    state Down
      do x = last x - v
      until not c then Up
  end
tel
\end{lstlisting}

The example above, from~\cite{heptreax_manual}, defines a component whose state variable \texttt{x}
(and thus output \texttt{y}) either increments or decrements based on the state of an automaton. The
contract enforces that the output \texttt{y} must remain within the range $[0, 10]$. The boolean
\texttt{c} is declared as a controllable variable using the \texttt{with} clause. The task of the
compiler is to synthesize a control strategy for \texttt{c}. To satisfy the contract, it must
generate logic that sets \texttt{c} to true (switching to the \texttt{Down} state) when \texttt{x}
approaches 10, and sets \texttt{c} to false (switching to the \texttt{Up} state) when \texttt{x}
approaches 0, effectively creating a bounded counter.

\paragraph*{Comparison with GRust}
Both \GRust and \Heptagon/BZR integrate formal specifications directly into the source code, but
their objectives differ fundamentally. \Heptagon/BZR focuses on \emph{synthesis}: automatically
generating an implementation from a high-level contract. In contrast, \GRust, through \GReusot,
focuses on the \emph{deductive verification} of a manually written implementation. This aligns with
a more traditional engineering workflow where the developer writes the control logic and a verifier
confirms its correctness.

\subsection{Verified Compilers and Code-Level Analyzers}
\label{sec:sota:verification:compilers}

Ensuring that the safety guarantees of a high-level model are preserved after compilation is a
critical challenge. A bug in the compiler could invalidate all high-level proofs. This is addressed
in two main ways: verifying the compiler itself, or verifying its output.

The most robust solution is a \emph{formally verified compiler}. An example is
CompCert~\cite{compcert}, a \Ci compiler formally proven in \Rocq to be free of miscompilation
errors. In the synchronous world, \Velus~\cite{velus} provides a formally verified compiler for a
subset of \Lustre, offering the strongest guarantee that the semantics of the source program are
preserved. However, the immense effort required means that evolving a language with a verified
compiler is extremely difficult, as even small changes can require significant proof maintenance.

An alternative approach is to analyze the generated code. This is the strategy taken by the \LustreC
language~\cite{lustrec}, which uses \FramaC~\cite{framac} for the formal verification of properties
directly on the generated \Ci code. Using techniques like weakest precondition calculus, \FramaC can
prove the absence of runtime errors (\eg buffer overflows) and verify functional contracts. While
this approach is powerful because it verifies the final artifact that will be deployed, it remains
extremely labor-intensive. A significant portion of the proof effort is dedicated to the onerous
task of establishing memory safety properties that are trivial in the high-level source model.

\paragraph*{Comparison with GRust}
\GRust offers a pragmatic and powerful alternative that sits between these two extremes. The \GRust
compiler itself is not formally verified like \Velus. However, its verification strategy directly
addresses the model-to-implementation gap in a novel way. The \GReusot tool translates contracts
written on the high-level \GRust model into formal annotations on the \emph{generated \Rust code}.
This means the verification effort targets the final artifact, ensuring the correctness of the
output of the compiler for a given program, without requiring a full proof of the compiler itself.

\subsection{Synthesis: GRust in the Verification Landscape}
\label{sec:sota:verification:synthesis}

The survey of verification techniques reveals a landscape of trade-offs between the level of
automation, the strength of the guarantees, and the practicality of the engineering workflow. The
approaches range from automated synthesis from contracts, to model-level verification, to laborious
code-level analysis. \GRust, through its verification tool \GReusot, carves out a distinct and
pragmatic position within this landscape, defined by a series of deliberate design choices.

First, \GReusot aligns with the contract-based deductive verification approach of tools like \Kind,
rather than the synthesis-from-contract model of \Heptagon/BZR. This choice supports a traditional
and flexible engineering workflow where the developer retains full control over the implementation
logic, and the formal method is used to verify its correctness, not to generate it.

Second, and most critically, \GReusot addresses the ``trust in tools" problem by verifying the
\emph{generated code}, not the abstract source model. This is philosophically similar to the
\LustreC/\FramaC toolchain, acknowledging that proofs on the final artifact are the most meaningful.
However, \GRust's choice of \Rust as a compilation target eliminates the generated assertions
required to prove memory safety in \Ci with \FramaC. The ownership and borrow-checking system of the
\Rust compiler provides these guarantees by construction, for free.

This allows the developer and the verifier (\Creusot) to focus exclusively on the high-level
functional properties of the system, as expressed in the contracts. \GReusot combines the rigor of
code-level analysis with the practicality of a modern, safe-by-construction language ecosystem,
offering a powerful and accessible path to high-assurance software.

\section{Overall Synthesis and Positioning of GRust}
\label{sec:sota:overall_synthesis}

This chapter has surveyed the state of the art across four key domains: industrial model-based
tools, academic synchronous languages, distributed systems frameworks, and Functional Reactive
Programming. The analysis reveals a spectrum of design philosophies, with different sets of
trade-offs between formal guarantees, performance, and programming abstraction.

Legacy tools like \Simulink and \Scade offer mature modeling environments but are tied to an
inefficient periodic execution model. Synchronous languages such as \Lustre and \Esterel provide
powerful formalisms for correctness, with derivatives like \Prelude and \textsc{HipHop.js}
adapting these principles to multi-periodic or asynchronous environments. \frp variants solve
theoretical issues of resource management, while Modal \frp introduces advanced type systems to
guarantee safety. Conversely, frameworks like ROS and \ReactiveX embrace asynchrony but delegate
the burden of ensuring correctness to the developer.

\GRust carves out a distinct and pragmatic position within this landscape, as summarized in
\Cref{tab:sota:synthesis:comparison}. It synthesizes principles from these different domains to
directly address the challenges of the Software-Defined Vehicle. It combines the declarative,
formally-grounded components of synchronous languages with the efficient, event-driven execution
model of practical reactive libraries, all while leveraging the safety guarantees of the \Rust
ecosystem.

\begin{table}[ht!]
  \centering
  \begin{tabularx}{\textwidth}{@{} >{\raggedright}p{0.15\textwidth}
    >{\raggedright}p{0.27\textwidth}
    >{\raggedright}p{0.27\textwidth}
    >{\raggedright\arraybackslash}p{0.27\textwidth} @{}}
    \toprule
    \textbf{Framework}                                                              & \textbf{Execution Model} & \textbf{Concurrency Guarantee} & \textbf{Primary Use Case} \\
    \midrule
    \Lustre / \Scade                                                                &
    Synchronous: pull-based execution at a single base-clock rate.                  &
    Sequential: compiles to a single-threaded step function.                        &
    Safety-critical control systems.                                                                                                                                        \\
    \addlinespace
    \Esterel                                                                        &
    Imperative synchronous: complex control-flow with preemption.                   &
    Deterministic concurrency: resolved at compile-time into a state machine.       &
    Time-critical control logic and hardware synthesis.                                                                                                                     \\
    \addlinespace
    \textsc{HipHop.js}                                                              &
    Synchronous reactive with asynchronous tasks.                                   &
    Deterministic concurrency: within the reactive machine.                         &
    Complex web applications.                                                                                                                                               \\
    \addlinespace
    \Prelude                                                                        &
    Multi-periodic, time-triggered: partitions a synchronous model into RTOS tasks. &
    Schedulability analysis: generates verifiable multi-task code.                  &
    Multi-periodic hard real-time systems.                                                                                                                                  \\
    \addlinespace
    \textsc{Lingua Franca}                                                          &
    Event-triggered with logical time: reactions ordered by timestamped events.     &
    Deterministic concurrency: enforced by a centralized runtime scheduler.         &
    Distributed real-time systems.                                                                                                                                          \\
    \addlinespace
    \Yampa                                                                          &
    Discrete-time, pull-based sampling of a continuous-time model.                  &
    Sequential: composition of stateful signal functions.                           &
    Animations, robotics, and GUIs.                                                                                                                                         \\
    \addlinespace
    \textsc{Async Rattus}                                                           &
    Asynchronous modal FRP: execution on local-clocks.                              &
    Causality and productivity by construction.                                     &
    GUIs and formal methods research.                                                                                                                                       \\
    \addlinespace
    ROS / \ReactiveX                                                                &
    Event-triggered: asynchronous message/event streams.                            &
    Non-deterministic: requires manual locking or schedulers.                       &
    Robotics and UI development.                                                                                                                                            \\
    \midrule
    \GRust                                                                          &
    GALS: locally synchronous components in a globally asynchronous system.         &
    Locally deterministic: guaranteed within services, but globally best-effort.    &
    Asynchronous automotive stack for SDV.                                                                                                                                  \\
    \bottomrule
  \end{tabularx}
  \caption[Comparative analysis of languages and frameworks.]%
  {A comparative analysis of languages and frameworks, positioning \GRust within the design spectrum
    of reactive and distributed systems.}
  \label{tab:sota:synthesis:comparison}
\end{table}
